% ===== 原书 p.3 =====
\section{指标问题}

\subsection{引言}

本文的中心议题是 Atiyah--Singer 指标定理。这是二十世纪后半叶最重要的数学成果之一，它在代数拓扑、微分几何与分析之间建立起根本的联系。具体地说，指标定理把线性椭圆微分算子的指标用若干拓扑不变量表示出来。寻找这种表示的问题最初由 Israel Gel'fand 于 1960 年提出，1962 年秋，Michael Atiyah 与 Isadore Singer 在哈佛与 MIT 解决了它。1963 年初，Atiyah 与 Singer 在《Bulletin of the American Mathematical Society》上联名发表论文 [9]，首次宣布指标定理，其中用配边论把指标表示成示性类的形式。1968 年，Atiyah 与 Singer 在《Annals of Mathematics》上以三篇论文的形式发表了指标定理的第二个证明（该系列的第二篇实际出自 Atiyah 与 Graeme Segal 之手）。第二个证明用 K 理论重新表述了指标定理，完全避开了上同调与配边论，并把它推广为等变的 $G$-指标定理与椭圆算子族的版本。1968 年以来，指标定理出现了多种不同的证明（例如热方程证明与超对称证明），它在数学物理中也变得重要起来。

在本文中，我们从第 1 节开始讨论 K 理论，这是指标定理的基本语言。K 理论主要由 Atiyah 与 Hirzebruch 在 1960 年代初发展起来；一般性的参考文献是 Atiyah 的经典著作 [6]。第 1 节末尾简要回顾示性类，以及 K 理论与上同调的关系。第 2、3 节沿 [10] 与 [11] 的思路给出指标定理的证明，其中第 2 节稍作停顿，介绍指标定理的上同调形式并给出几个简单应用。计入紧 Lie 群 $G$ 的作用之后的等变推广放在第 4 节。等变 K 理论由 Atiyah 与 Segal 在牛津发展起来，完整的处理见 Segal 的 [17]。在第 4 节中，我们先陈述 $G$-指标定理，然后按 Raoul Bott 的论文 [12] 考察齐性微分算子这一特殊情形。第一个附录利用一族微分算子的指标给出 Bott 周期性定理的一个证明；这部分讨论基本沿 Atiyah 的论文 [5]，并把该技巧推广，用来证明 Thom 同构定理的通常形式与等变形式。第二个附录不假定任何同调代数知识，直接导出 Atiyah--Hirzebruch 谱序列。

书中的材料并非原创；事实上，所有结果在作者出生之前就已为人所知。本书要做的是把 1963 至 1968 年这一关键时期 K 理论与指标理论各方面的发展综合起来，重点放在 Atiyah 及其各位合作者的工作上。书中试图统一这一领域各篇论文的记号，并详细算出若干例子。特别是，第 2 节对奇维流形上指标的处理，以及附录中关于粘合函数的例子，都比现有文献详细得多。

本书是为高年级本科生与研究生写的指标定理导论，假定他们几乎没有或完全没有 K 理论与指标理论的基础。不过，书中假定读者熟悉向量丛的基本理论；若不熟悉，应先浏览 [16] 的前几节或 [6] 的第 1 节，再读本书。建议读者具备较好的上同调论基础并了解一些示性类，但对指标定理本身的证明而言，这些并非严格必需。第 4 节还需要一点表示论的知识，可参阅 [1]。

% ===== 原书 p.4 =====
\subsection{定义与例子}

设 $X$ 是紧光滑流形。给定 $X$ 上的光滑复向量丛 $E$ 与 $F$，我们分别用 $\Gamma E$ 与 $\Gamma F$ 记 $E$ 与 $F$ 的 $C^\infty$ 截面构成的空间（确切地说，$\Gamma V$ 由满足 $\pi\circ s=\id$ 的光滑映射 $s:X\to V$ 组成，其中 $\pi:V\to X$ 是向量丛 $V$ 的投影映射）。

\begin{definition}
映射 $P:\Gamma E\to\Gamma F$ 称为\emph{线性偏微分算子}，如果它在 $X$ 的局部坐标 $x_1,\dots,x_n$ 下取形式
\[
P=\sum_{r\le m}\alpha_{i_1\cdots i_r}\,
\frac{\partial^{i_1+\cdots+i_r}}{\partial x_{i_1}\cdots\partial x_{i_r}},
\]
其中 $\alpha_{i_1\cdots i_r}(x):E_x\to F_x$ 是 $E$ 与 $F$ 的纤维之间的线性变换，光滑依赖于 $x$。
\end{definition}

换言之，$P$ 局部地是算子 $\partial/\partial x_j$ 的多项式，系数取矩阵值并光滑依赖于 $x$。形式地把未定元 $i\xi_j$ 代入 $\partial/\partial x_j$，我们得到相应的多项式 $p(x,\xi)$，其中 $\xi=(\xi_1,\dots,\xi_n)$。为说明多项式与算子之间的联系，可把算子写成 $P=p(x,D)$ 的形式，其中 $D=(-i\partial/\partial x_1,\dots,-i\partial/\partial x_n)$。\footnote{读者可能认出，映射 $\partial/\partial x_j\mapsto i\xi_j$ 就是 Fourier 变换。这一与 Fourier 分析的联系在本文后面会变得重要。} 若只保留多项式 $p(x,\xi)$ 中次数最高的那些项，我们得到齐次多项式，记作 $\sigma(x,\xi)$，称为 $P$ 的\emph{象征}。

\begin{definition}
线性偏微分算子称为\emph{椭圆的}，如果它的象征 $\sigma(x,\xi)$ 对一切 $x\in X$ 与一切非零向量 $\xi=(\xi_1,\dots,\xi_n)\neq0\in\R^n$ 都可逆。
\end{definition}

\begin{example}
设 $X=\R^n$，$E=F=\R^n\times\C$ 是平凡复线丛。此时 $\Gamma E$ 或 $\Gamma F$ 的元素就是 $\R^n$ 上的光滑复值函数。$\R^n$ 上的 Laplace 算子是二阶微分算子
\[
\nabla^2=\left(\frac{\partial}{\partial x_1}\right)^2+\cdots+\left(\frac{\partial}{\partial x_n}\right)^2 .
\]
它的象征是 $\sigma(x,\xi)=-\left((\xi_1)^2+\cdots+(\xi_n)^2\right)<0$（$\xi\neq0$），所以 Laplace 算子是椭圆的。更一般地，若取形如 $D=f\nabla^2$ 的算子，其中 $f$ 是光滑的非零复值函数，则 $D$ 也是椭圆算子。
\end{example}

\begin{example}
令 $X=\C=\R^2$，取 $z=x+iy$，其中 $x$、$y$ 是实坐标。复分析中的 Cauchy--Riemann 算子是
\[
\frac{\partial}{\partial\bar z}=\frac12\left(\frac{\partial}{\partial x}
+i\frac{\partial}{\partial y}\right).
\]
回忆 Cauchy--Riemann 微分方程由 $\partial f/\partial\bar z=0$ 给出。$\partial/\partial\bar z$ 的象征是 $\sigma(x,\xi)=\frac12 i(\xi_1+i\xi_2)$；由于 $\xi$ 是实向量，只要 $\xi\neq0$ 就有 $\sigma\neq0$，所以 $\partial/\partial\bar z$ 是椭圆算子。我们将在第 2 节末尾回到这个算子。
\end{example}

给定紧流形 $X$ 上的椭圆算子 $D:\Gamma E\to\Gamma F$，偏微分算子理论告诉我们，$D$ 的核（即 $Df=0$ 的全部解构成的空间）与 $D$ 的余核（即商空间 $\Gamma F/\operatorname{im}D$）都是有限维的。于是可以定义

% ===== 原书 p.5 =====
\begin{definition}
紧流形上椭圆算子 $D$ 的\emph{指标}由
\[
\Ind D=\dim\Ker D-\dim\Coker D
\]
给出。
\end{definition}

\begin{example}
设 $X$ 是圆周 $S^1=\R/2\pi\Z$，考虑算子 $d/d\theta$。它的象征是函数 $\sigma(x,\xi)=i\xi$，所以该算子是椭圆的。核是微分方程 $\frac{d}{d\theta}f(\theta)=0$ 的全部解构成的空间，显然就是常函数。为计算余核，注意 $S^1$ 上任何光滑函数 $f(\theta)$ 都可写成 Fourier 级数 $f(\theta)=\sum_{-\infty}^{\infty}a_ne^{in\theta}$。于是得到
\[
\frac{d}{d\theta}f(\theta)=\sum_{-\infty}^{\infty}in\,a_ne^{in\theta}.
\]
注意该级数中常数项（$n=0$）消失了，可见余核同样只是常函数。因此
\[
\Ind d/d\theta=1-1=0.
\]
\end{example}

更一般地，考虑由 $P=p(d/d\theta)$ 给出的椭圆算子，其中 $p(\xi)$ 是常系数多项式。用 Fourier 级数展开，立即可得
\[
Pf(\theta)=\sum_{-\infty}^{\infty}p(in)\,a_ne^{in\theta}.
\]
设 $n_1,\dots,n_r$ 是方程 $p(in)=0$ 的互不相同的整数解，则 $P$ 的核由一切形如 $f(\theta)=b_1e^{in_1\theta}+\cdots+b_re^{in_r\theta}$ 的函数组成，故 $\dim\Ker P=r$。又注意 $P$ 的像由一切满足 $b_{n_1}=\cdots=b_{n_r}=0$ 的函数 $f(\theta)=\sum_{-\infty}^{\infty}b_ne^{in\theta}$ 组成，于是可以把 $P$ 的余核与 $P$ 的核等同起来，故 $\dim\Coker P=r$。由此得到 $\Ind P=r-r=0$。事实上，本文后面会看到，圆周 $S^1$ 上任何线性椭圆偏微分算子的指标都是零。

按我们的定义，椭圆算子的指标是一个与某些偏微分方程的解相联系的整数。本文的目标是把指标用更容易计算的拓扑不变量表示出来，这些不变量涉及向量丛、上同调群与 Chern 类。按代数拓扑的精神，我们处理这一问题的方法可由下图很好地概括：
\[
\begin{tikzcd}[row sep=2.8em, column sep=2.8em]
K(TX) \arrow[rr, "\mathrm{t\text{-}ind}"] & & \Z \\
& \text{椭圆算子} \arrow[ul, "\sigma"'] \arrow[ur] &
\end{tikzcd}
\]

\transnote{图中三个映射分别在后面三处定义：$\sigma$ 是第 2.1 节的象征映射，$\mathrm{t\text{-}ind}$ 是第 2.2 节构造的拓扑指标，而未标出名字的对角箭头就是第 0.2 节定义的解析指标 $\Ind$（第 3.3 节）。第 2、3 节的全部内容可以概括为一句话：这个三角形交换。指标定理的证明并不是去算两边的具体数值，而是验证两条路径给出同一个整数。}


接下来三节将逐一解释图中的各个方面。我们已经看到从椭圆算子到整数的指标映射的定义。第 1 节将定义并研究函数 $K(\cdot)$ 的性质以及它与上同调的关系。在第 2 节中，我们引入 K 理论中的象征映射 $\sigma$ 与拓扑指标 $\mathrm{t\text{-}ind}$。Atiyah--Singer 指标定理的实际证明将在第 3 节给出，其内容就是证明上图交换。在这一过程中，必须把讨论扩展到椭圆伪微分算子类。

% ===== 原书 p.6 =====
\section{K 理论与上同调}

\subsection{向量丛与 $K(X)$}

设 $X$ 是紧 Hausdorff 空间，$\Vect(X)$ 是 $X$ 上复向量丛的同构类构成的空间。Whitney 和 $\oplus$ 使 $\Vect(X)$ 成为一个带单位元的阿贝尔半群。于是可以按 Grothendieck 的形式作差手续，构造相应的虚拟向量丛群 $K(X)$。

\begin{definition}
$K(X)$ 是 $\Vect(X)\times\Vect(X)$ 中 $\Delta\Vect(X)$ 的陪集构成的空间，其中 $\Delta:\Vect(X)\to\Vect(X)\times\Vect(X)$ 是对角映射。换言之，我们把 $K(X)$ 定义为向量丛的一切对 $E-F$，模去等价关系
\[
E^1-E^2=F^1-F^2
\iff
\exists\,G \text{ 使得 } E^1\oplus F^2\oplus G\cong E^2\oplus F^1\oplus G .
\]
这样得到的集合 $K(X)$ 是阿贝尔群，负元由 $-(E-F)=(F-E)$ 给出。
\end{definition}

也可以把 $K(X)$ 定义为以 $\Vect(X)$ 为生成元的自由阿贝尔群，模去由一切形如 $E+F-(E\oplus F)$ 的元素生成的子群。两种定义殊途同归。向量丛的张量积 $\otimes$ 可以延拓为 $K(X)$ 上的乘法，从而赋予 $K(X)$ 带单位元的交换环结构。

给定 $X$ 上的向量丛 $E$，用 $[E]$ 记它在 $K(X)$ 中的像。注意映射 $\Vect(X)\to K(X)$ 不必是单射。例如考虑球面 $S^2\subset\R^3$ 上的实向量丛：法丛 $NS^2$ 同构于平凡线丛 $\mathbf{1}$，于是在 $K_\R(S^2)$ 中得到
\[
[TS^2]=[3]-[NS^2]=[3]-[1]=[2],
\]
而切丛 $TS^2$ 与平凡平面丛 $[2]$ 并不同构。这正是必须在上述等价关系中引入 $G$ 项的原因。因此，若 $[E]=[F]$，则存在向量丛 $G$ 使 $E\oplus G\cong F\oplus G$；再取向量丛 $G^\perp$ 使 $G\oplus G^\perp\cong\mathbf{n}$ 对某个 $\mathbf{n}$ 成立。由此可见，$[E]=[F]$ 当且仅当 $E$ 与 $F$ 稳定等价。

\transnote{等价关系中的 $G$ 依赖于被比较的两对丛，而不是事先固定的。若把 $G$ 固定下来，「存在 $G$ 使 $E\oplus G\cong F\oplus G$」这种关系就没有传递性，无法定义陪集。取过 $G$ 后再补上正交补，一切就归约为稳定等价：$K(X)$ 的元素不必是某个具体丛的类，而是「虚拟丛」。}


给定光滑映射 $f:X\to Y$ 与 $Y$ 上的向量丛 $F$，可以构造 $X$ 上的拉回丛 $f^*F$，它对一切 $x\in X$ 满足 $(f^*F)_x=F_{f(x)}$。这一映射进而延拓为环同态 $f^*:K(Y)\to K(X)$，称为 $f$ 的诱导映射。由此可见，$K(\cdot)$ 是从紧空间范畴到带单位元交换环范畴的反变函子。特别地，若 $i:\pt\to X$ 是点到 $X$ 的包含映射，则得到映射 $i^*:K(X)\to K(\pt)\cong\Z$。\footnote{点上的向量丛就是向量空间，它由维数决定，在同构意义下唯一。} 定义\emph{约化 K 群}为 $\tilde K(X)=\Ker i^*$，就得到分裂
\[
K(X)\cong\tilde K(X)\oplus K(\pt).
\]
把 $\tilde K(X)$ 看作 $K(X)$ 模去平凡丛所得的商，则 $\tilde K(X)$ 由稳定向量丛的类组成。

\transnote{$i^*$ 就是把向量丛映为它在一点处的纤维维数，即取秩。因此 $\Ker i^*$ 由正负两部分秩相等的虚拟丛组成；$\tilde K(\pt)=0$，而 $K(\pt)=\Z$。实 K 理论中 $K_\R(\pt)=\Z$，复 K 理论中同样是 $\Z$，故上面这个分裂把「秩」这一平凡信息干净地分离出去了。}


\subsection{复形与紧支撑}

现在考虑 $X$ 不紧、而只是局部紧（且 Hausdorff）的情形。设 $X^+$ 是 $X$ 的单点紧化（若 $X$ 本身已紧，就定义 $X^+=X\cup\pt$），则可以讨论紧支撑的 K 理论。

\begin{definition}
若 $X$ 局部紧，则定义
\[
K(X):=\tilde K(X^+)=K(X^+)/K(\pt).
\]
\end{definition}

注意，若 $X$ 本身已紧，这一定义与前面的定义一致。把讨论限制在逆紧映射上，\footnote{\emph{逆紧映射}是指紧集的原像仍为紧集的映射。这一性质使逆紧映射可以延拓为两个空间的单点紧化之间的连续映射。} 这一定义就把 $K(\cdot)$ 延拓为从局部紧 Hausdorff 空间范畴到交换环范畴的反变函数（注意：$K(\cdot)$ 是带单位元的环，当且仅当 $X$ 紧）。\footnote{用紧支撑的观点，$K(\cdot)$ 也可以看作关于开集包含的共变函数。设 $i:U\to X$ 是开子集 $U$ 到 $X$ 的包含映射，定义\emph{自然延拓同态} $i_*:K(U)\to K(X)$ 为由 $X^+\to X^+/(X^+-U^+)\cong U^+$ 诱导的映射。}

% ===== 原书 p.7 =====
换一个角度，也可以考虑 $X$ 上向量丛的复形 $E^\bullet$，其形式为
\[
0\to E^0\xrightarrow{\alpha_0}E^1\xrightarrow{\alpha_1}\cdots
\xrightarrow{\alpha_{n-1}}E^n\to0,
\]
其中 $\alpha_k\alpha_{k-1}=0$。$X$ 上的两个复形 $E^\bullet$ 与 $F^\bullet$ 称为\emph{同伦的}，若存在 $X\times[0,1]$ 上的复形 $G^\bullet$，使 $E^\bullet$ 与 $F^\bullet$ 分别同构于限制 $G^\bullet|_{X\times0}$ 与 $G^\bullet|_{X\times1}$。该复形的\emph{支撑}是 $X$ 中使复形限制到 $x$ 处纤维上不成为正合列的一切点 $x$ 构成的集合。我们关心的是具有紧支撑的复形。

\begin{claim}
设 $C(X)$ 表示 $X$ 上具有紧支撑的复形按紧支撑同伦构成的同伦类集合，$C_\emptyset(X)$ 是 $C(X)$ 中由空支撑复形构成的子集，则
\[
K(X)\cong C(X)/C_\emptyset(X),
\]
其中 $K(X)$ 是上面定义的紧支撑 K 群。
\end{claim}

证明见 [6, 第 2.6 节] 或 [17, 附录]。若 $X$ 实际上是紧的，则复形 $E^\bullet$ 对应于元素
\[
\chi(E^\bullet)=\sum_k(-1)^k[E^k]\in K(X).
\]
这个交错和映射 $\chi:C(X)\to K(X)$ 就是通常的\emph{Euler 特征}。

除了任意长度的复形，也可以只用固定长度 $n>0$ 的复形，仍然得到 $K(X)\cong C^n(X)/C^n_\emptyset(X)$。特别地，把复形限制为长度 $1$ 往往很方便，即「在紧集外同构的向量丛对 $(E,F)$」。设 $\varphi:C^1(X)\to C(X)$ 是显然的包含映射，我们如下构造它的左逆 $\psi:C(X)\to C^1(X)$。给定任意长度的复形
\[
0\to E^0\xrightarrow{\alpha_0}E^1\xrightarrow{\alpha_1}\cdots
\xrightarrow{\alpha_{n-1}}E^n\to0,
\]
在每个向量丛 $E^k$ 上取 Hermite 内积 $\langle\cdot,\cdot\rangle_k$。于是得到伴随映射 $\alpha_k^*:E^{k+1}\to E^k$，由性质
\[
\langle\alpha_k(v),w\rangle_{k+1}=\langle v,\alpha_k^*(w)\rangle_k
\]
刻画。现在令 $\psi(E^\bullet)$ 为相应的长度 $1$ 复形
\[
0\to\bigoplus_i E^{2i}\xrightarrow{\alpha}\bigoplus_i E^{2i+1}\to0,
\]
其中
\[
\alpha=\bigoplus_i\alpha_{2i}+\bigoplus_i\alpha_{2i+1}^*.
\]
容易验证 $E^\bullet$ 与 $\psi(E^\bullet)$ 有相同的支撑。另外注意，复形 $\psi(E^\bullet)$ 与内积的选取无关，因为所有内积彼此同伦。

\transnote{把偶数项直接并置、奇数项先取伴随再并置，是 Hodge 理论中「$d+d^*$」这一手法的离散版本：伴随项负责把上同调「压回」核与余核。正是这个符号安排使 $\psi$ 成为左逆，从而任意长度的复形与长度 $1$ 的复形给出同一个 $K$ 群；若漏掉伴随，$\alpha$ 一般不再是复形的微分。}


% ===== 原书 p.8 =====
采用任意长度复形的好处是，乘法的定义变得方便。给定 $X$ 上的向量丛 $E$ 与 $Y$ 上的向量丛 $F$，定义它们在 $X\times Y$ 上的\emph{外张量积}为
\[
E\BT F:=\pi_X^*E\otimes\pi_Y^*F,
\]
其中 $\pi_X$ 与 $\pi_Y$ 分别是 $X\times Y$ 到 $X$ 与 $Y$ 的坐标投影。取复形的外张量积，即可把它延拓为乘积映射
\[
\BT:K(X)\otimes K(Y)\to K(X\times Y),
\]
其中 $(E^\bullet\BT F^\bullet)^i=\bigoplus_jE^j\BT F^{i-j}$，且
\[
\alpha_i^{E^\bullet\BT F^\bullet}
=\bigoplus_j\Big(\big(\alpha_j^{E^\bullet}\BT1_{F^{i-j}}\big)
+(-1)^j\big(1_{E^j}\BT\alpha_{i-j}^{F^\bullet}\big)\Big).
\]

例如，给定长度 $1$ 的复形
\[
0\to E^0\xrightarrow{\alpha}E^1\to0
\qquad\text{与}\qquad
0\to F^0\xrightarrow{\beta}F^1\to0,
\]
它们分别定义在 $X$ 与 $Y$ 上，其外张量积是 $X\times Y$ 上的复形
\[
0\to E^0\BT F^0
\xrightarrow{\alpha\BT1+1\BT\beta}
E^1\BT F^0\oplus E^0\BT F^1
\xrightarrow{-1\BT\beta+\alpha\BT1}
E^1\BT F^1\to0 .
\]
注意必须引入因子 $-1$，才能使 $\alpha_1\alpha_0=0$。按上述构造，相应的长度 $1$ 复形为
\[
0\to E^0\BT F^0\oplus E^1\BT F^1
\xrightarrow{\theta}
E^0\BT F^1\oplus E^1\BT F^0\to0,
\]
其中
\[
\theta=
\begin{pmatrix}
\alpha\BT1 & -1\BT\beta^*\\
1\BT\beta & \alpha^*\BT1
\end{pmatrix}.
\]

\subsection{齐次复形}

我们引入紧支撑 K 理论的主要动机，是考虑 $V$ 上的复形，其中 $V$ 是 $X$ 上的实向量丛，投影映射为 $\pi:V\to X$（后文将取 $V$ 为 $TX$，即 $X$ 的余切丛）。此时可以在复形上附加齐次结构。给定 $X$ 上的复向量丛 $E$ 与 $F$，可把它们提升为 $V$ 上的向量丛 $\pi^*E$ 与 $\pi^*F$，并且有 $(\pi^*E)_v=E_{\pi(v)}$ 与 $(\pi^*F)_v=F_{\pi(v)}$。若对一切 $v\in V$ 与实数 $\lambda>0$ 都有
\[
\alpha_{\lambda v}=\lambda^m\alpha_v:E_{\pi(v)}\to F_{\pi(v)},
\]
则称同态 $\alpha:\pi^*E\to\pi^*F$ 是（正）$m$ 次\emph{齐次的}。

给定 $V$ 上的度量，可见 $\alpha$ 完全由它在球丛 $S(V)$ 上的限制决定。设 $X$ 紧，$E^\bullet$ 是 $V$ 上形如
\[
0\to\pi^*E^0\xrightarrow{\alpha}\pi^*E^1\xrightarrow{\alpha}\cdots
\xrightarrow{\alpha}\pi^*E^n\to0
\]
的复形，其中 $\alpha^2=0$，且 $\alpha$ 是 $m$ 次齐次的。若 $E^\bullet$ 在 $S(V)$ 上正合，则 $E^\bullet$ 的支撑就是零截面，即 $X$ 在 $V$ 中的像。于是 $E^\bullet$ 具有紧支撑，从而代表 $K(V)$ 的一个元素。\footnote{注意，若 $m<0$，则 $\alpha$ 在零截面上不连续。所幸我们不需要 $m<0$ 的情形。不过，若 $X$ 只是局部紧的，就必须取 $m=0$，此时 $\alpha$ 在零截面上可能不连续。见 [10, 第 492--493 页]。} 我们断言，$K(V)$ 可以用这类齐次复形来定义。

\transnote{原文用 $TX$ 记 $X$ 的\emph{余切}丛，而通行记号 $TX$ 表示切丛，余切丛写作 $T^*X$。本译本保留原文字面，因此下文凡出现 $K(TX)$，都应理解为余切丛的 K 群。这样记也有好处：象征 $\sigma(x,\xi)$ 是余向量 $\xi$ 的齐次多项式，其取值自然落在余切丛的 K 群中。}

% ===== 原书 p.9 =====
\begin{claim}
记 ${}^{m}C(V)$ 为 $V$ 上 $m$ 次齐次复形的同伦类集合，${}^{m}C_\emptyset(V)$ 为其中使同态 $\alpha$ 沿球丛 $S(V)$ 的每个纤维为常值的那部分正合复形。则
\[
{}^{m}C(V)/{}^{m}C_\emptyset(V)\cong C(V)/C_\emptyset(V)\cong K(V).
\]
\footnote{与前述相同，可以把讨论限制在固定长度的齐次复形上。}
\end{claim}

\begin{proof}
考虑包含 ${}^{m}C(V)\subset C(V)$，我们想构造一个逆映射 $C(V)\to{}^{m}C(V)$。设 $E^\bullet\in C(V)$ 是支撑为紧集 $L$ 的复形
\[
0\longrightarrow E^0\xrightarrow{\ \alpha_0\ }E^1\xrightarrow{\ \alpha_1\ }\cdots\xrightarrow{\ \alpha_{n-1}\ }E^n\longrightarrow 0,
\]
在 $V$ 上取一个度量，并令 $D(V)$ 是包含 $L$ 的圆盘丛。把 $E^i$ 限制到零截面，得 $F^i=E^i|_X$，于是在 $D(V)$ 上有 $E^i\cong\pi^*F^i$。定义 $\beta_i:\pi^*F^i\to\pi^*F^{i+1}$ 为在相应的球丛 $S(V)$ 上与 $\alpha_i$ 一致的那个 $m$ 次齐次映射，这样就得到 $V$ 上的一个齐次复形 $F^\bullet$：
\[
0\longrightarrow\pi^*F^0\xrightarrow{\ \beta_0\ }\pi^*F^1\xrightarrow{\ \beta_1\ }\cdots\xrightarrow{\ \beta_{n-1}\ }\pi^*F^n\longrightarrow 0.
\]
容易看出 $E^\bullet$ 与 $F^\bullet$ 同伦，因此它们对应 $K(V)$ 中的同一个元素；于是映射 $C(V)\to{}^{m}C(V)$ 是同构。其次，若 $E^\bullet\in C_\emptyset(V)$ 的支撑为空，则 $E^\bullet$ 同伦于一个使 $\alpha_i$ 沿 $V$ 的每个纤维为常值的复形，由此可见 $F^\bullet\in{}^{m}C_\emptyset(V)$。反之，若 $F^\bullet\in{}^{m}C_\emptyset(V)$ 使 $\beta_i$ 沿 $S(V)$ 的每个纤维为常值，则取同伦 $\alpha_i(v,t)=\|v\|^{tm}\beta_i(v/\|v\|)$，可见 $F^\bullet$ 同伦于某个正合复形 $E^\bullet\in C_\emptyset(V)$。于是 ${}^{m}C(V)/{}^{m}C_\emptyset(V)\cong C(V)/C_\emptyset(V)$。
\end{proof}
